\scanpage{23}
% lines 1–3
un intérêt (contrairement aux formes de $\mathrm{SL}(2)_S$ ou $\mathrm{GP}(1)_S$, qui sont équivalentes et équivalent aux formes de $\mathbb P^1_S\ldots$).

% lines 4–6
Revenant au cas général, il peut être intéressant de calculer la forme \deleted{\inserted{quadratique}} de Casimir de $\mathfrak g=E'$ et $\mathfrak s=E$, \unsure{définis|définies} par

% lines 7
\[
\deleted{\operatorname{Cas}_{E}(\unsure{\eta',\eta'|x,y})=\rho_{E'}\operatorname{Tr}\rho_E}
\tag*{\deleted{(73)}}
\]

% lines 8–9
\[
\deleted{\left\{\begin{aligned}
\operatorname{Cas}_{\unsure{\mathfrak g'|E'}}(x')\deleted{\unsure{,y'}})&=\operatorname{Tr}\rho_{E'}(x')^2\\
\operatorname{Cas}_E(x)&=\operatorname{Tr}\widetilde\rho_E(x)^2
\end{aligned}\right.}
\tag*{\deleted{(73)}}
\]

% lines 10–11
\deleted{(les formes bil. associées sont respectivement $2\operatorname{Tr}\rho_{E'}(x')\rho_{E'}(y')$ et $2\operatorname{Tr}\widetilde\rho_E(x)\widetilde\rho_E(y)$).}

% lines 12
\deleted{Le calcul donne :}

% lines 13
\[
\deleted{\operatorname{Cas}_E(x^{\unsure{'}})=\Delta\varphi\overset{\deleted{\unsure{x',x'|x,x}}}{(x,x)}} \tag*{\deleted{(74)}}
\]

% lines 14
\deleted{d'où, en remplaçant $(E,\omega,\varphi)$ par $(E',\omega',\varphi')$}

% lines 15
\[
\deleted{\operatorname{Cas}_E(x)=} \tag*{\deleted{(75)}}
\]

% lines 16–17
\[
\left\{\begin{aligned}
\operatorname{Cas}_{E'}(x',y')&=\operatorname{Tr}\bigl(\rho_{E'}(x')\mathbin{.}\rho_{E'}\unsure{(}y')\bigr)\\
\operatorname{Cas}_{E}(x,y)&=\operatorname{Tr}\bigl(\widetilde\rho_E(x)\mathbin{.}\widetilde\rho_E(y)\bigr)
\end{aligned}\right. \tag*{(73)}
\]

% lines 18
On trouve alors

% lines 19–20
\[
\left\{\begin{aligned}
\operatorname{Cas}_{E'}(x',y')&=2\varphi'(x',y')\\
\operatorname{Cas}_E(x,y)&=-2\Delta\varphi(x,y)
\end{aligned}\right. \tag*{(74)}
\]

% lines 21–23
Pour le voir, on note d'abord que les \inserted{deux} formules ont la propriété d'homogénéité qu'il faut, quand on multiplie $\omega$ par $\lambda$, $\varphi$ par $\mu$ : les

% lines foot
% archive sheet number: 22
